\scanpage{301}
- \mspage{2} -
Tribu 43

invariance du nombre d'éléments de la p-base. Rabiots sur $F^{p^k}$ (cf. bas de la p.22). Ensuite un chapeau disant que, comme toute dérivation de L sur K est nulle sur $K(L^p)$ on est réduit à étudier les dérivations des ext. radicielles (``du premier étage''). La prop.2 ; un cor. disant qu'il existe une dérivation envoyant n'importe où les éléments d'une famille p-indépendante ; donner aussi une condition nécessaire et suffisante pour que les $dx_i$ engendrent.

Le n$^\circ$ 4 (p.23-24) des applications. La prop.4 peut aller en exer. La prop.5 est un bon énoncé ; y nommer a) et b) les deux phrases. . Dans la prop.6, il faut préciser qu'il s'agit de dérivations dans un corps contenant F ; l'énoncer sous forme d'équivalence de : séparabilité, - les dérivations se prolongent, - les différentielles s'injectent.

Il faudra ajouter le sorite à la Lang sur les dimensions d'espaces de dérivations (``le sous truc, la tour, la translation et le compositum'')

\uline{Fabrication d'extensions}.

Ça doit venir à la fin du \S\ sur les extensions cycliques. On exposera la méthode générale (cf. papier que vendra Weil), l'exemple de la théorie de Kummer, celui d'Artin-Schreier simple (cyclique d$^\circ$p), la wittification entre astérisques, et peut être ce qu'on pourra extraire de Klein ou du bouquin de Tchébotarof sur le groupe icosaédrique.

II. - \uline{Matières devant aller en} Alg. \uline{Commutative}, \uline{chapitre des anots affines}.

Une idée du contenu général de ce chapitre est donnée dans les décisions générales. Les choses suivantes proviennent des ``Suppléments'' :

\uline{Extensions régulières et consorts}.

Il s'agit des n$^\circ$2,3,4, du \S1, plus le lemme 2 du n$^\circ$1. Le Congrès objecte au traitement parallèle des ext. régulières, primaires et séparables

\mspage{300}
